\subsection{Team 2}
\label{app:team2}

\paragraph{Data} The challenge training set (\chk{2,992} scans with findings and
masks), with the \chk{200}-scan validation set used only for model selection.
Annotations were used unmodified. No external or in-house CT data, pseudo-labels
or derived anatomy masks were used in the final system. The only externally
pretrained component was Qwen3-Embedding-4B \citep{zhang2025qwen3} (Apache 2.0),
used frozen, and the image network was trained from scratch. Test-set scans and
finding texts were not used for training or model selection. Whether any such
material was part of the text encoder's own pretraining data is not known.

\paragraph{Pre-processing} Following nnU-Net preprocessing \citep{isensee2021nnu},
volumes were oriented to LPS and resampled to about
\chk{0.7 $\times$ 0.7 $\times$ 1.0}\,mm. Intensities were clipped to
$[-1000, 197]$\,HU and normalized with a mean of $-549$\,HU and a standard
deviation of 313\,HU. Training used random foreground-centred crops with zero
padding, with patch sizes from \chk{224 $\times$ 224 $\times$ 128} to
\chk{448 $\times$ 448 $\times$ 256} voxels.

\paragraph{Architecture and text encoder} The model, Grounder, uses a 3D nnU-Net
residual-encoder (ResEnc-L) encoder \citep{isensee2024nnu} and a five-stage
decoder. The frozen Qwen3-Embedding-4B converts each finding into a
2,560-dimensional embedding, and two-layer MLPs generate scale and shift
parameters that condition the decoder features through FiLM
\citep{perez2018film}, $F' = \gamma F + \beta$.

\paragraph{Training} Dice plus cross-entropy loss, AdamW with a fixed learning
rate of $10^{-4}$, and a batch size of one. Models generally converged within
20--30 epochs. Image augmentation included rotation, scaling and intensity
transformations, but no flipping. In-distribution text augmentation replaced a
finding with a real finding from another training case that matched it in
location and finding type, keeping the original CT and mask. Each model trained
on one H200 GPU for about 4--7 days, depending on patch size.

\paragraph{Ensembling} Thirteen independently trained models differing in patch
size, selected on validation performance, with voxel-wise probabilities averaged
equally.

\paragraph{Post-processing} Probabilities were thresholded at 0.5. No
connected-component filtering, anatomical restriction, or rewriting of the text
at inference was applied.

\paragraph{Inference time and hardware} About 5\,s per scan per finding on one
H200 GPU.

\paragraph{Final submission and lessons} The final submission was
Grounder\_1.4n (full test set: Dice \chk{0.365}, hit rate \chk{0.763}, instance
F1 \chk{0.343}). Larger anatomical context, in-distribution text augmentation
and diverse patch sizes improved grounding. Adding lobe maps from
TotalSegmentator \citep{wasserthal2023totalsegmentator} or category embeddings
did not help, and augmentation with negative text was unreliable because the
training annotations are partial.

\paragraph{Code and weights} The repository \url{https://github.com/khuhm/Grounder}
(MIT licence) will hold the code and trained weights, which have not yet been
released.
